% ============================================================
\chapter{Retrieval-Augmented Generation}
\label{ch:rag}
% ============================================================

Large language models are remarkably capable, but they have two structural weaknesses: they ``hallucinate'' (generate plausible but false information) and their knowledge is frozen at training time. In 2020, Patrick Lewis and collaborators at Facebook AI Research proposed an elegant solution: \emph{Retrieval-Augmented Generation} (RAG). The principle is to couple the generative model with a document retrieval system: before generating a response, the model retrieves relevant passages from an external knowledge base and uses them as context. This architecture decouples \emph{memory} (the document base, easily updated) from \emph{reasoning} (the language model). RAG is today at the heart of enterprise question-answering systems, specialized chatbots, and research assistants.

\begin{intuition}[title=\textbf{Why augment with retrieval?}]
Generative language models (LLMs) suffer from hallucinations and
knowledge frozen at training time. Retrieval-Augmented Generation (RAG)
addresses these problems by providing the model with relevant documents
from an external knowledge base before generation. This grounds the
output in verifiable facts and enables knowledge updates without
retraining.
\end{intuition}

% ------------------------------------------------------------
\section{RAG Framework}
\label{sec:rag:framework}
% ------------------------------------------------------------

\begin{definition}[Retrieval-Augmented Generation]
\index{RAG}
A \textbf{RAG} system combines a \emph{retriever} and a \emph{generator}.
Given a query $q$ and a corpus $\mathcal{C}$:
\begin{enumerate}
    \item \textbf{Retrieve}: extract the $k$ most relevant documents
    $\{d_1, \ldots, d_k\} = \text{Retrieve}(q, \mathcal{C})$
    \item \textbf{Generate}: produce the answer
    $y = \text{Generate}(q, d_1, \ldots, d_k)$
\end{enumerate}
\end{definition}

\begin{keyformulas}[title=\textbf{Probabilistic RAG Formulation}]
\[
    P(y | q) = \sum_{d \in \mathcal{C}} P(d | q) \, P(y | q, d)
    \approx \sum_{i=1}^{k} P(d_i | q) \, P_{\text{LLM}}(y | q, d_i)
\]
where $P(d | q)$ is the relevance score and $P_{\text{LLM}}(y | q, d)$
is the generative probability conditioned on the document.
\end{keyformulas}

\begin{proposition}[Benefits of RAG]
\begin{enumerate}
    \item \textbf{Up-to-date knowledge}: the knowledge base can be updated
    without retraining the model.
    \item \textbf{Traceability}: sources are identifiable (citations).
    \item \textbf{Reduced hallucinations}: the model grounds its output
    in verifiable facts.
    \item \textbf{Parameter efficiency}: a smaller model with retrieval
    can match the quality of a larger model.
\end{enumerate}
\end{proposition}

% ------------------------------------------------------------
\section{Dense Retrieval}
\label{sec:rag:dense}
% ------------------------------------------------------------

\begin{definition}[Dense Passage Retrieval (DPR)]
\index{DPR}\index{dense retrieval}
DPR (Karpukhin et al., 2020) encodes queries and documents into a
shared vector space:
\[
    \text{sim}(q, d) = \mathbf{e}_q^\top \mathbf{e}_d
    = E_Q(q)^\top E_D(d)
\]
where $E_Q$ and $E_D$ are BERT encoders fine-tuned to maximize
similarity between queries and relevant passages.
\end{definition}

\begin{definition}[Contrastive Training Loss for DPR]
Training uses a contrastive loss over positive and negative passages:
\[
    \mathcal{L} = -\log \frac{\exp(\text{sim}(q, d^+))}
    {\exp(\text{sim}(q, d^+)) + \sum_{j=1}^{n}
    \exp(\text{sim}(q, d_j^-))}
\]
where $d^+$ is the relevant passage and $\{d_j^-\}$ are negatives
(in-batch negatives and hard negatives from BM25).
\end{definition}

\begin{definition}[ColBERT]
\index{ColBERT}
ColBERT (Khattab et al., 2020) uses a \emph{late interaction} between
query and document tokens:
\[
    \text{sim}(q, d) = \sum_{i=1}^{\abs{q}} \max_{j=1}^{\abs{d}}
    \mathbf{q}_i^\top \mathbf{d}_j
\]
where $\mathbf{q}_i$ and $\mathbf{d}_j$ are contextual token embeddings.
This fine-grained interaction improves precision while remaining
efficient through precomputed document embeddings.
\end{definition}

\begin{remark}
Dense retrieval outperforms traditional sparse methods (BM25) on many
benchmarks, especially for semantic matching. However, BM25 remains
competitive for exact keyword matching and is often used as a first-stage
retriever or for generating hard negatives.
\end{remark}

% ------------------------------------------------------------
\section{Vector Databases}
\label{sec:rag:vectordb}
% ------------------------------------------------------------

\begin{definition}[Vector Database]
\index{vector database}
A \textbf{vector database} stores embedding vectors and supports
Approximate Nearest Neighbor (ANN) search in sub-linear time. Examples
include FAISS, Pinecone, Weaviate, Qdrant, and ChromaDB.
\end{definition}

\begin{proposition}[Indexing Methods]
The main ANN indexing methods are:
\begin{itemize}
    \item \textbf{IVF} (Inverted File): partitions the space into
    Voronoi cells for coarse search.
    \item \textbf{HNSW} (Hierarchical Navigable Small World): a navigable
    graph structure offering excellent speed-accuracy trade-offs.
    \item \textbf{PQ} (Product Quantization): compresses vectors via
    quantization to reduce memory footprint.
\end{itemize}
\end{proposition}

% ------------------------------------------------------------
\section{Chunking Strategies}
\label{sec:rag:chunking}
% ------------------------------------------------------------

\begin{definition}[Document Chunking]
\index{chunking}
\textbf{Chunking} divides documents into fragments suitable for the
model's context window:
\begin{itemize}
    \item \textbf{Fixed-size}: $n$ tokens with overlap of $m$ tokens.
    \item \textbf{Paragraph-based}: respects document structure.
    \item \textbf{Semantic}: splits at topic boundaries (using a
    segmentation model).
    \item \textbf{Recursive}: hierarchical splitting
    (title $\to$ section $\to$ paragraph).
\end{itemize}
\end{definition}

\begin{warning}[title=\textbf{Chunk Size Trade-off}]
Chunks that are too small lose context; chunks that are too large dilute
relevant information and waste the model's token budget. In practice,
$256$--$512$ tokens with $50$--$100$ overlap offers a good balance.
\end{warning}

% ------------------------------------------------------------
\section{RAG Architectures}
\label{sec:rag:architectures}
% ------------------------------------------------------------

\begin{definition}[RAG-Sequence and RAG-Token]
\index{RAG-Sequence}\index{RAG-Token}
Lewis et al.\ (2020) propose two variants:
\begin{itemize}
    \item \textbf{RAG-Sequence}: a single document is used for the entire
    generation:
    $P(y|q) \approx \sum_{d} P(d|q) \prod_t P(y_t | q, d, y_{<t})$
    \item \textbf{RAG-Token}: the document can change at each token:
    $P(y|q) \approx \prod_t \sum_{d} P(d|q) P(y_t | q, d, y_{<t})$
\end{itemize}
\end{definition}

\subsection{Advanced RAG}

\begin{definition}[Self-RAG]
\index{Self-RAG}
Self-RAG (Asai et al., 2023) teaches the model to decide \emph{when}
to retrieve (via special reflection tokens) and to evaluate the
relevance of retrieved documents before using them.
\end{definition}

\begin{definition}[RAPTOR]
\index{RAPTOR}
RAPTOR builds a hierarchical tree of summaries: leaves are original
chunks, internal nodes are cluster summaries. Retrieval traverses the
tree from general to specific.
\end{definition}

\begin{intuition}
Modern RAG pipelines are increasingly sophisticated: they include
reranking of retrieved documents, query rewriting, iterative retrieval
(the model can issue follow-up queries), and self-reflection on whether
the retrieved context is sufficient.
\end{intuition}

% ------------------------------------------------------------
\section{Evaluating RAG Systems}
\label{sec:rag:eval}
% ------------------------------------------------------------

\begin{definition}[RAG Evaluation Metrics]
\index{RAG evaluation}
RAG evaluation covers three dimensions:
\begin{enumerate}
    \item \textbf{Retrieval quality}: recall@$k$, precision@$k$, MRR
    (Mean Reciprocal Rank).
    \item \textbf{Faithfulness}: is the answer consistent with the
    retrieved documents?
    \item \textbf{Answer relevance}: does the answer correctly address
    the question?
\end{enumerate}
\end{definition}

\begin{keyformulas}[title=\textbf{Recall@$k$ and MRR}]
\[
    \text{Recall@}k = \frac{\abs{\{d \in \text{top-}k : d \in \mathcal{R}\}}}
    {\abs{\mathcal{R}}}
    \qquad
    \text{MRR} = \frac{1}{\abs{Q}} \sum_{q \in Q} \frac{1}{\text{rank}(d_q^+)}
\]
where $\mathcal{R}$ is the set of relevant documents and $d_q^+$ is the
first relevant document.
\end{keyformulas}

\begin{example}[RAGAS Framework]
The RAGAS framework automatically evaluates RAG systems on four metrics:
faithfulness, answer relevance, context relevance, and context recall.
It uses an LLM as a judge to assess each dimension.
\end{example}

% ------------------------------------------------------------
\section{Python Implementation}
\label{sec:rag:code}
% ------------------------------------------------------------

\begin{codeblock}[title=\textbf{Minimal RAG Pipeline}]
\begin{minted}{python}
from sentence_transformers import SentenceTransformer
import numpy as np

class SimpleRAG:
    def __init__(self, docs, model_name="all-MiniLM-L6-v2"):
        self.encoder = SentenceTransformer(model_name)
        self.docs = docs
        self.embeddings = self.encoder.encode(docs)

    def retrieve(self, query, k=3):
        q_emb = self.encoder.encode([query])
        scores = (self.embeddings @ q_emb.T).squeeze()
        top_k = np.argsort(scores)[-k:][::-1]
        return [(self.docs[i], scores[i]) for i in top_k]

    def generate_prompt(self, query, k=3):
        results = self.retrieve(query, k)
        context = "\n\n".join([doc for doc, _ in results])
        return (f"Context:\n{context}\n\n"
                f"Question: {query}\nAnswer:")

# Usage
docs = ["Paris is the capital of France.",
        "Berlin is the capital of Germany.",
        "The Eiffel Tower is 330 meters tall."]
rag = SimpleRAG(docs)
prompt = rag.generate_prompt("What is the capital of France?")
\end{minted}
\end{codeblock}

% ------------------------------------------------------------
\section{Exercises}
\label{sec:rag:exercises}
% ------------------------------------------------------------

\begin{exercise}[RAG Pipeline]
Implement a complete RAG pipeline using ChromaDB as the vector store and
a HuggingFace model for generation. Test on a corpus of Wikipedia
articles.
\end{exercise}

\begin{exercise}[Chunking Strategies]
Compare three chunking strategies (fixed-size, paragraph-based, semantic)
on a corpus of 100 documents. Measure recall@5 for a set of 50 questions.
\end{exercise}

\begin{exercise}[DPR vs BM25]
Compare DPR (dense retrieval) with BM25 (sparse retrieval) on the
Natural Questions dataset. In which cases does dense retrieval outperform
BM25? And vice versa?
\end{exercise}

\begin{exercise}[RAG Evaluation]
Design an evaluation protocol for a medical question-answering RAG
system. Which metrics are critical? How should hallucinations be handled
in a medical context?
\end{exercise}
